\chapter{Resultados}
\label{ch:resultados}

Este capítulo apresenta os resultados das duas frentes de coleta descritas no Capítulo~\ref{ch:metodologia}: as entrevistas com gestores dos quatro estabelecimentos, complementadas pela entrevista com um fornecedor de cardápios digitais, e o questionário com frequentadores. Ao final, as duas fontes são integradas conforme o plano de análise da Seção~\ref{sec:analise}.

\section{Casos estudados}
\label{sec:res-casos}

\pendente{caracterizar os casos A, B e C (perfil, bairro, ano de abertura, porte, modalidade de autoatendimento) e resumir cada entrevista após a transcrição}

\subsection{Caso D, Ordinária Pub}
\label{sec:res-caso-d}

\pendente{confirmar o ano de abertura do bar}

O Caso D é o Ordinária Pub, bar no bairro Bom Fim aberto depois de 2020, cuja identificação no trabalho foi autorizada pelo entrevistado. O entrevistado foi um dos sócios, que havia deixado a sociedade cerca de dois meses antes da entrevista; o bar segue em funcionamento, e por isso ele descreve a operação no passado. A entrevista ocorreu em 13 de novembro de 2026, por videoconferência, e teve cerca de 18 minutos transcritos. Segundo ele, o público típico é jovem adulto, de 20 a 40 e poucos anos, da própria região, ampliado pelos shows de samba, que atraíam pessoas de outras zonas da cidade e de faixas etárias mais altas.

\textbf{Atendimento e cardápio.} O atendimento era feito na mesa, por garçons, apenas com cardápio físico. O único uso digital do cardápio era a publicação da mesma arte nos destaques do Instagram do bar, como divulgação.

\textbf{Motivos da escolha.} O entrevistado atribui a decisão principalmente a uma preferência pessoal dos sócios, que não gostam de cardápio digital como consumidores. A escolha não foi motivada por custo nem por dificuldade técnica: o sistema de gestão contratado pelo bar oferecia um módulo de cardápio digital sem custo adicional, o que tornaria a implantação simples. Os motivos citados foram:

\begin{itemize}
    \item \textbf{Navegação.} No papel, é possível ver a página inteira e comparar itens rapidamente; no celular, a rolagem dificulta voltar e comparar opções. Nas palavras do entrevistado, ``no papel não tem scroll''.
    \item \textbf{Manutenção e desgaste.} QR Codes colados nas mesas perdem a tinta com a limpeza, e tablets e totens quebram, o que torna o custo de manutenção parecido com o de reimprimir cardápios, com a diferença de que o digital facilita atualizar preços e itens.
    \item \textbf{Contato humano.} O atendimento pelo garçom permite tirar dúvidas e receber sugestões, o que ele considera parte da experiência de um bar.
\end{itemize}

O entrevistado reconhece vantagens do cardápio digital, como atualizar preços automaticamente e retirar na hora itens que acabaram, evitando o constrangimento de o garçom avisar que um item do cardápio está indisponível. Sua principal crítica não é ao digital em si, mas à falta de alternativa: ``O que me incomoda muito é quando não tem [\ldots] a segunda opção'', porque isso cria uma barreira para quem tem pouco letramento digital, está sem internet ou não tem smartphone. Ele também observou que ``todo passo que tu adiciona, ele é uma barreira e é um ponto de ruído''.

\textbf{Tecnologia nos bastidores.} Embora o cliente não usasse nenhuma interface digital, a operação era fortemente informatizada. O bar usava o sistema de gestão Saipos, contratado como serviço: os garçons lançavam os pedidos no sistema, que os enviava às impressoras, processava o pagamento integrado ao banco e gerava relatórios financeiros. O estoque era atualizado pela importação das notas fiscais de compra e baixado automaticamente a cada pedido, a partir das fichas técnicas associadas aos itens do cardápio, com alertas de estoque mínimo que alimentavam uma lista de compras.

\textbf{Visão de futuro.} Para o entrevistado, uma solução digital só seria preferível ao papel se permitisse ver muitos itens ao mesmo tempo, como uma folha A4 ou A5; ele mencionou como exemplo hipotético um cardápio projetado na mesa. Tablets e totens, que já encontrou como consumidor, repetiriam o problema da rolagem e acrescentariam custo e risco de quebra para o estabelecimento.

\section{Fornecedor de cardápio digital}
\label{sec:res-fornecedor}

\pendente{caracterizar a empresa e resumir a entrevista, com foco em arquitetura, requisitos, dados pessoais e efeito da lei municipal}

\section{Questionário com frequentadores}
\label{sec:res-questionario}

\subsection{Participação e perfil da amostra}
\label{sec:res-perfil}

\pendente{número total de respostas; respostas excluídas por falta de consentimento; tamanho da base principal (Porto Alegre e região metropolitana) e da base de outras localidades; número de respondentes que nunca encontraram cardápios digitais}

% Tabela sugerida: perfil da base principal (faixa etária, gênero, escolaridade,
% frequência de ida, tipo de estabelecimento, deficiência), com n e %.

\pendente{tabela e comentário sobre o perfil da amostra, comparando com o público esperado e registrando o viés da amostragem por conveniência}

\subsection{Experiência e preferência}
\label{sec:res-experiencia}

\pendente{frequência com que os respondentes encontram cardápios digitais e preferência geral (físico, digital, indiferente)}

\subsection{Construtos UTAUT2 e extensões contextuais}
\label{sec:res-construtos}

% Tabela sugerida: por construto, média, mediana, desvio-padrão, % de
% concordância (4 e 5) e correlação entre os dois itens, após recodificar
% EE2, HM2 e PV2. Extensões (TRU, PRI, SOC) item a item.

\pendente{tabela de escores por construto e comentário sobre os construtos mais e menos favoráveis}

\subsection{Problemas técnicos e situações de uso}
\label{sec:res-problemas}

\pendente{frequência de cada problema técnico relatado, taxa de desistência ou frustração e categorias das respostas abertas (funciona melhor, funciona pior, relatos de frustração)}

\subsection{Cruzamentos exploratórios}
\label{sec:res-cruzamentos}

\pendente{escores e preferência por faixa etária, frequência de ida, deficiência ou dificuldade e localidade (Porto Alegre e região metropolitana versus outras cidades)}

\section{Integração das fontes}
\label{sec:res-integracao}

\pendente{para cada categoria das entrevistas, indicar se o questionário traz evidência convergente, divergente ou complementar}
